\section{The unsmoothed collision spectrum on a bounded frequency set}
\label{sec:compact-collision-spectrum}
The flight-age inversion reduces the original physical observation to
three collision frequencies. We now estimate these frequencies without
smoothing the collision record. The relevant multiplier is the record
of the \emph{preceding} flight, on the image side of the transfer
operator. Its discontinuities are transverse to the stable test curves.
For a bounded set of real roof frequencies the action identity controls
the accumulated multiplier on every homogeneous inverse branch. This
allows the collision-space compactness argument to be applied to the
actual displacement--time twist. No induced operator is introduced.

Put
\[
 f_R^{\rm c}=(\kappa_{R,1},\kappa_{R,2},\tau_R-\bar\tau_R),\qquad
 \check f_R^{\rm c}=f_R^{\rm c}\circ T_R^{-1},\qquad
 Q_{R,z}=M_{\exp(i z\cdot\check f_R^{\rm c})}\mathcal L_R.
\]
The last expression will first be defined on smooth densities and then
on the collision spaces. Real frequencies are denoted by
$\omega=(u,b)\in\T^2\times\R$. Near zero, $z\in\mathbb C^3$ is
an ordinary coordinate vector. We use the covariance $\Sigma_R$ in
\eqref{eq:physical-three-collision-record}; equivalently it is the
collision Green--Kubo covariance of $f_R^{\rm c}$.

\subsection{The image-side multiplier}
\begin{lemma}[A bounded analytic realization of the exact collision twist]
\label{lem:backward-flight-multiplier}
The local common collision spaces can be chosen so that $Q_{R,z}$ is
analytic in $z$ in the strong operator norm. On each compact complex
frequency set its norm and all fixed-order derivatives are bounded
uniformly in $R$. For smooth functions $a,d$ and every $m\ge0$,
\begin{equation}\label{eq:exact-collision-endpoint-pairing}
 \ell\bigl(M_d Q_{R,z}^{m}(a\nu)\bigr)
 =\int_M a(x)\exp\!\left(i z\cdot\sum_{j=0}^{m-1}
                  f_R^{\rm c}(T_R^j x)\right)d(T_R^m x)\,d\nu(x).
\end{equation}
In particular the earliest density occurs on the right of the operator
product and the terminal multiplier on the left.
\end{lemma}
\begin{proof}
We specify the admissible discontinuities. In the collision coordinates
$(\alpha,\varphi)$, the velocity obtained by following the incoming
flight backwards is
$w=\cos\varphi\,n_\alpha-\sin\varphi\,t_\alpha$. Its length to a
candidate preceding disk with center $d$ is the entry root
\[
 (d-Rn_\alpha)\cdot w-
 \sqrt{R^2-|d-Rn_\alpha|^2+((d-Rn_\alpha)\cdot w)^2}.
\]
Finite horizon gives a fixed finite list of candidates. On a region
where one regular candidate is selected, this root is uniformly
$1/2$-H\"older, up to its boundary. Indeed the quantities outside and
under the square root are uniformly Lipschitz and bounded, and
$|\sqrt x-\sqrt y|\le\sqrt{|x-y|}$ for $x,y\ge0$.
The preceding displacement is the constant label $-d$ on that region.
The assertion is made in $(\alpha,\varphi)$, not by an inverse change
of variables at grazing.

The boundaries between preceding itineraries belong to the one-step
backward singularity set. They are unstable curves, uniformly
transverse to the stable cone. Splitting into their finitely many
monotone pieces gives a partition with a uniformly bounded number of
components in each homogeneity strip. Each stable curve meets a
partition element in a uniformly bounded number of intervals; the
length within distance $\epsilon$ of its boundary is $O(\epsilon)$.
These are the two partition conditions of \cite[Lemma 3.3]{DPZ}.
They concern the preceding label, not the forward label.

Choose the weak stable-test exponent $p$ with $0<p<1/6$, a strong
stable exponent $0<q<p$, and an unstable comparison exponent
$0<\gamma<\min\{p-q,1/12\}$, decreasing it also to satisfy the
geometric restrictions of the spaces. Choose the multiplier exponent
$\zeta$ strictly between $\max\{p,\gamma/(1-\gamma)\}$ and $1/2$.
The strict inequality at $1/2$ places the preceding observables in the
closure used in the H\"older test spaces. The piecewise multiplier
bound of \cite[Lemma 3.3(b)]{DPZ} therefore applies. This is the same
image-side placement as in its Lemma 3.9, now including the entry-root
coordinate. The exponential series and all its fixed derivatives
converge in this multiplier norm on compact complex sets. This proves
boundedness and analyticity.

For one step, the convention $\mathcal L h(\psi)=h(\psi\circ T)$ gives
\[
 Q_{R,z}h(\psi)
 =h\bigl(e^{iz\cdot f_R^{\rm c}}\psi\circ T_R\bigr).
\]
Iteration proves \eqref{eq:exact-collision-endpoint-pairing} on
smooth densities. Continuity gives the asserted realization. The local
boundary reparametrizations and the deck-invariant restriction in
Lemma~\ref{lem:collision-spectral-input} preserve these bounds.
\end{proof}

\subsection{Accumulated phases in the collision norms}
We record the geometric estimate needed for real frequencies away from
zero. It is not a consequence of analyticity near zero.

\begin{lemma}[Action bounds on homogeneous and matched curves]
\label{lem:collision-action-weights}
For $|b|\le B$, let
\[
 w_{m,\omega}(x)=\exp\{i(u\cdot K^{\rm c}_{m,R}(x)
                         +bS_{m,R}(x)-bm\bar\tau_R)\}.
\]
On every homogeneous component of $T_R^{-m}W$, $W\in\mathcal W^s$,
this weight has modulus one and uniformly bounded Lipschitz seminorm,
with bound $C_B$ independent of $m,R,u$. For the matched inverse pieces
used in the unstable comparison of the collision norm, parametrized
over the same interval and arising from terminal curves at distance
at most $\epsilon$, their two weights satisfy
\begin{equation}\label{eq:matched-action-weight}
 \|w_{m,\omega}^{(1)}-w_{m,\omega}^{(2)}\|_{C^q}
                       \le C_B\epsilon^{1-q}.
\end{equation}
The unmatched pieces are exactly the geometric unmatched pieces; the
weight creates no additional partition boundary.
\end{lemma}
\begin{proof}
Write $\vartheta_R=R\sin\varphi\,d\alpha$. The one-flight action
identity, proved in Section~\ref{sec:measurable-phase-rigidity}, is
$d\tau_R=T_R^*\vartheta_R-\vartheta_R$. On any connected arc
$\gamma$ regular through $m$ collisions it telescopes to
\begin{equation}\label{eq:long-action-connector}
 S_{m,R}(y)-S_{m,R}(x)
       =\int_{T_R^m\gamma}\vartheta_R-\int_\gamma\vartheta_R.
\end{equation}
Here $x,y$ are the endpoints of $\gamma$; the lattice sum is constant
on the corresponding regular branch. The form is uniformly bounded.
A homogeneous inverse stable component is a graph with bounded slope,
and its image under $T_R^m$ contracts by at most $C\Lambda^{-m}$ in
arc length. The latter is the stable-curve estimate used in
\cite[equation (4.6)]{DZ}. Apply
\eqref{eq:long-action-connector} to each subarc. Its two lengths are
bounded by a constant times the distance of its endpoints in the graph
coordinate. This proves the first assertion, uniformly in $m$.

We give the matching check because a bound on the variation along only
one curve would not suffice. In the construction of
\cite[Section 4.3]{DZ}, the untrimmed matched inverse pieces are joined
by unstable-cone connectors $\gamma_x$ such that
\[
 |\gamma_x|\le C\Lambda^{-m}\epsilon,
              \qquad |T_R^m\gamma_x|\le C\epsilon.
\]
Every intermediate image of such a connector is regular; otherwise the
piece is assigned to the unmatched part. The two endpoints consequently
have identical lattice itineraries. Formula
\eqref{eq:long-action-connector} bounds their roof-sum difference by
$C\epsilon$. Trimming to a common graph interval moves the paired
point along one inverse stable piece by at most $C\epsilon$; its image
moves by at most $C\Lambda^{-m}\epsilon$. The first part of the proof
bounds this additional change also by $C\epsilon$. Thus the two weights
in common graph coordinates differ in supremum norm by $C_B\epsilon$.
Their individual Lipschitz constants in those coordinates are bounded
by $C_B$. Interpolation between the supremum and Lipschitz bounds gives
\eqref{eq:matched-action-weight}. All connectors and all trimming here
are those of the existing collision partition, so no new singular
boundary or long-word count is introduced.
\end{proof}

\begin{proposition}[Bounded-band Lasota--Yorke estimates]
\label{prop:bounded-band-LY}
For every $B<\infty$, on a finite cover by local common collision
spaces there are $C_B<\infty$, $0<\theta<1<\Lambda$, and $C_3\ge1$
such that, for all real $(u,b)\in\T^2\times[-B,B]$ and $m\ge0$,
\begin{align}
 |Q_{R,\omega}^m h|_w&\le C_B|h|_w,
                                               \label{eq:roof-LY-weak}\\
 \|Q_{R,\omega}^m h\|_s
 &\le C_B(\theta^m+\Lambda^{-qm})\|h\|_s+C_B|h|_w,
                                               \label{eq:roof-LY-stable}\\
 \|Q_{R,\omega}^m h\|_u
 &\le C_B\Lambda^{-\gamma m}\|h\|_u+C_B C_3^m\|h\|_s.
                                               \label{eq:roof-LY-unstable}
\end{align}
An equivalent strong norm, chosen after $B$, makes this a uniformly
power-bounded, quasi-compact family with essential spectral radius
strictly less than one. On each local parameter neighborhood,
$Q_{R,\omega}$ is continuous from the strong space to the weak space
as $(R,\omega)$ varies, uniformly on this frequency set.
\end{proposition}
\begin{proof}
We check the changes to the norm estimates rather than importing a
local limit theorem for a different record. The collision norms and
the inverse-curve decomposition are those in \cite[Section 3.3]{DZ};
its Proposition 4.1 gives the unweighted estimates with precisely the
three forms above. A density paired against a terminal stable test
$\psi$ has, on each inverse piece $U$, the integral
\begin{equation}\label{eq:weighted-curve-pairing}
 \int_U h\,w_{m,\omega}\,J_U T_R^m\,
                                  (\psi\circ T_R^m)\,dm_U.
\end{equation}
The collision probability is invariant, so there is no exponential
volume-Jacobian factor. We use the same density convention as the
unweighted norms.

For the weak estimate, Lemma~\ref{lem:collision-action-weights} bounds
the $C^p$ norm of the added weight uniformly. The pulled-back test and
the normalized stable Jacobian have the same $C^p$ bounds as before.
The sum of the stable Jacobian suprema over inverse pieces is bounded
uniformly in $m$ \cite[Lemma 3.2(b)]{DZ}. This proves
\eqref{eq:roof-LY-weak}.

For the strong stable estimate, on $U$ subtract the average of
$\psi\circ T_R^m$, \emph{not} the average of the whole weighted test.
The $C^q$ norm of the resulting difference is at most
$C\Lambda^{-qm}|\psi|_{C^q}$, by stable contraction. Multiplying this
difference by the uniformly Lipschitz weight keeps that decaying factor.
The average term times the weight and normalized Jacobian is a $C^p$
test and is paid by the weak norm on pieces which have grown to the
fixed reference length. Pieces not yet reaching that length retain the
usual geometric growth-lemma factor $\theta^m\|h\|_s$. The associated
Jacobian sums, including their length normalization, are unchanged,
because $|w_{m,\omega}|=1$. These are exactly the two sums in
\cite[equations (4.7)--(4.11)]{DZ}. This proves
\eqref{eq:roof-LY-stable}; a nondecaying derivative of the phase has
been put in the weak term, not in the leading strong coefficient.

For the unstable estimate, retain the matched and unmatched pieces of
that proof. Unmatched pieces use the same stable norm estimate, since
the extra weight has bounded $C^q$ norm. On matched pieces split the
test difference into the old Jacobian/test difference multiplied by one
weight and the difference of the weights multiplied by the second
Jacobian/test. The first term gives the old leading factor
$C_B\Lambda^{-\gamma m}\|h\|_u$ and its stable-norm error. For the
second term \eqref{eq:matched-action-weight}, divided by
$\epsilon^\gamma$, is bounded, since $1-q>\gamma$. The Jacobian and
length sums bound it by $C_B C_3^m\|h\|_s$, as for the stable-norm
error in \cite[equation (4.3)]{DZ}. We retain this exponential cross
coefficient; it is not asserted uniform in $m$.

Choose an integer $N$ so that the coefficients of $\|h\|_s$ and
$\|h\|_u$ in their respective first terms are less than $1/4$.
Then choose $c_B>0$ with $c_B C_B C_3^N<1/4$ and use
$\|h\|_B=\|h\|_s+c_B\|h\|_u$. It follows that
\[
 \|Q_{R,\omega}^N h\|_B\le\tfrac12\|h\|_B+C_B'|h|_w.
\]
Iterate this inequality using \eqref{eq:roof-LY-weak}, and treat the
finitely many remaining powers. This proves uniform power boundedness.
The compact embedding of the strong unit ball into the weak space,
\cite[Lemma 3.2(d)]{DPZ}, gives quasi-compactness with a uniform
essential-radius bound on the stated set. The equivalent norm does not
change eigenvalues or spectral projections.

Finally apply the one-step matched-curve comparison in
\cite[Section 5 and Theorem 2.6]{DZ} to two nearby radii. On unmatched
pieces the old small-length estimate is multiplied by a bounded
$C^p$ norm. On a matched pair the preceding disk label is the same;
the explicit preceding-flight root is uniformly $1/2$-H\"older in
$(R,\alpha,\varphi)$. Its difference tends uniformly to zero in
supremum norm, and interpolation with its uniform $C^{1/2}$ bound
makes the difference tend to zero in $C^q$. Thus the added weighted
matched term tends to zero in the strong-to-weak operator norm. The
same assertion in frequency follows from Lemma~\ref{lem:backward-flight-multiplier}.
The prescribed boundary-length reparametrizations are already part of
this comparison. This proves the asserted local uniform continuity.
\end{proof}

\subsection{Physical exclusion of the peripheral spectrum}
\begin{theorem}[The full compact-frequency collision estimate]
\label{thm:compact-collision-spectrum}
For every $B<\infty$ and $\epsilon>0$ there are $C_{B,\epsilon}<\infty$
and $\rho_{B,\epsilon}<1$ such that
\begin{equation}\label{eq:full-compact-frequency-gap}
 \|Q_{R,(u,b)}^m\|\le C_{B,\epsilon}\rho_{B,\epsilon}^m,
 \quad |b|\le B,\quad |(u,b)|\ge\epsilon,
 \quad u\in[-\pi,\pi]^2.
\end{equation}
For a fixed neighborhood of zero, uniform in $R$, the simple
eigenvalue and projection satisfy
\begin{align}
 Q_{R,z}^m&=\lambda_R(z)^m\Pi_R(z)+N_R(z)^m,
                 &\|N_R(z)^m\|&\le C\rho^m,
                                               \label{eq:raw-collision-splitting}\\
 \log\lambda_R(z)&=-\tfrac12 z^{\mathsf T}\Sigma_R z+O(|z|^3),
                 &\|\Pi_R(z)-\Pi_R(0)\|&\le C|z|.
                                               \label{eq:raw-collision-expansion}
\end{align}
Here $\Pi_R(0)h=\ell(h)\nu$. For real $z$ in a smaller neighborhood,
$|\lambda_R(z)|\le e^{-c|z|^2}$ with $c>0$ uniform in $R$.
\end{theorem}
\begin{proof}
Fix a real frequency. Quasi-compactness and power boundedness from
Proposition~\ref{prop:bounded-band-LY} leave only finite-dimensional,
semisimple eigenspaces on the unit circle. We show that a vector in
such an eigenspace is an actual measurable phase. If $E_\lambda$ is
its spectral projection, then on the strong space
\[
 E_\lambda=\lim_{M\to\infty}\frac1M
                       \sum_{j=0}^{M-1}\lambda^{-j}Q_{R,\omega}^j.
\]
Smooth densities are dense. Their images under $E_\lambda$ form a dense
linear subspace of the finite-dimensional range, hence its whole range.
Thus any vector there is $E_\lambda(a\nu)$ for some smooth $a$.
For every smooth test $\psi$, the exact pairing and invariance give
\[
 |Q_{R,\omega}^j(a\nu)(\psi)|
       \le\|a\|_\infty\int|\psi\circ T_R^j|\,d\nu
       =\|a\|_\infty\int|\psi|\,d\nu.
\]
The same bound holds in the Cesaro limit. The vector is therefore
$q\nu$ for a bounded measurable $q$. Invertibility and the eigenvector
equation imply that $|q|$ is invariant; ergodicity makes it constant.
For a nonzero vector normalize it to one. Then
\[
 q\circ T_R=\lambda^{-1}
   e^{i(u\cdot\kappa_R+b\tau_R-b\bar\tau_R)}q.
\]
Apply Theorem~\ref{thm:complete-physical-phase}, with constant phase
$-\arg\lambda-b\bar\tau_R$. It gives $u=0$ on the torus, $b=0$,
and $\lambda=1$. This argument uses the already proved physical phase
rigidity, not a choice of representative at isolated periodic points.

For every nonzero frequency the spectral radius is consequently less
than one. The uniform Lasota--Yorke estimates and strong-to-weak
continuity permit the spectral stability theorem
\cite[Theorem 2.1]{DZ} to be applied locally in $(R,u,b)$. A finite
cover of the compact set in \eqref{eq:full-compact-frequency-gap}
gives the displayed complementary power estimate. Compactness is used
only \emph{after} the common-space inequalities and the physical
peripheral exclusion have been proved.

Near zero, strong analyticity and the uniform unperturbed collision gap
give \eqref{eq:raw-collision-splitting}. The first derivative of the
eigenvalue vanishes because $\int f_R^{\rm c}\,d\nu=0$. Twice
differentiating \eqref{eq:exact-collision-endpoint-pairing} with
$a=d=1$, dividing by $m$, and using the spectral splitting identifies
the Hessian of $\log\lambda_R$ as minus the limiting collision
covariance. The exponentially summable covariance formula of
Theorem~\ref{thm:collision-covariance}, applied to $f_R^{\rm c}=P_Rh_R$,
identifies it with $\Sigma_R$. Cauchy's estimate on a fixed smaller
complex neighborhood gives the uniform cubic remainder. Its strict
positive lower eigenvalue bound then proves the final assertion.
\end{proof}

\begin{remark}
The constants in \eqref{eq:full-compact-frequency-gap} may depend on
$B$ without any specified growth bound. We will first let $m$ tend to
infinity on each fixed frequency set and then enlarge that set by
positive approximation. Substituting a growing $B_m$ directly into
\eqref{eq:full-compact-frequency-gap} is not part of the proof. The
operator also contains no section indicator and proves no corresponding
spectral assertion for the unbounded four-coordinate return twist.
\end{remark}
